\ifdefined\included
\else
\documentclass[a4paper,11pt,twoside]{StyleThese}
\include{formatAndDefs}
% fallback si compilation standalone
\providecommand{\acs}[1]{#1}
\providecommand{\minitoc}{}
\providecommand{\dominitoc}{}
\providecommand{\faketableofcontents}{}
\hbadness=10000
\sloppy
\begin{document}
\setcounter{chapter}{1} %% Numéro du chapitre précédent ;)
\dominitoc
\faketableofcontents
\fi

\chapter{Etat de l'art des attaques 5G}
\minitoc

% (PB1 : Manque d’information de la part des opérateurs sur ce qui est réaliste ou non)
Les opérateurs de télécommunications communiquent très peu sur la nature de leurs implémentations ainsi que sur les attaques auxquelles ils sont confrontés. Cette « sécurité par l’obscurité » limite l’émergence de menaces opportunistes et préserve des avantages concurrentiels. Toutefois, elle constitue également un frein important pour la recherche académique. En l’absence de données empiriques et de retours d’expérience concrets, de nombreux travaux proposent des scénarios d’attaque théoriques sans disposer de cadre de référence solide, rendant difficile l’évaluation de leur réalisme et de leur applicabilité dans des environnements opérationnels.

% (PB2 : La définition d’une attaque 5G est flou à cause de mauvaise caractérisation)
Ce flou de l'état de l'art en matière de réalisme est accentué par une caractérisation parfois insuffisante des attaques dans la littérature scientifique. Certains travaux reposent sur des modèles d’attaquants incomplets, des hypothèses peu explicites ou encore des expérimentations menées sur des implémentations open source trop peu matures. Bien qu’il existe des cadres méthodologiques dédiés à la caractérisation des attaques tels que FIGHT, ceux-ci présentent encore des limites et demeurent, à notre connaissance, très peu adoptés dans l’état de l’art. Cette situation contribue à entretenir une certaine ambiguïté quant à la définition et à la classification des attaques réellement pertinentes dans le contexte de la 5G.

% (PB3 : Proposition d’une nouvelle méthode de classification innovante)
Par ailleurs, les approches traditionnelles de détection d’anomalies reposent majoritairement sur l’analyse de flux à grande échelle, fondée sur des indicateurs volumétriques. Cependant, la softwarisation des réseaux 5G rapproche ces infrastructures de systèmes distribués complexes, ce qui nécessite désormais une analyse plus fine, notamment au niveau des paquets. Afin de répondre à ce changement de paradigme, nous proposons une nouvelle catégorisation en trois classes d'attaques : sémantiques, séquentielles et topologiques, chacune présentant des problématiques spécifiques en matière de détection.

\section{Les plateformes open source comme une approximation du réel}

\section{Exemples de mécanismes de sécurité négligés par l’état de l’art}

\section{Distinction entre les attaques réalistes et les hypothèses irréalistes }
\subsection{Attaques spécifiques aux implémentations}
\subsection{Attaques visant des failles dans les standards}

\section{Les limites de MITRE en tant que framework de catégorisation}
\subsection{Absence de techniques bien connues}
\subsection{Status incorrects des techniques}
\subsection{Manque de détails sur les caractéristiques des techniques}
\subsection{Pourquoi et comment caractériser nos attaqus}

\section{Nouvelle classification propre aux enjeux de détections}
\subsection{Anomalies sémantiques}
\subsection{Anomalies séquentielles}
\subsection{Anomalies topologiques}

\section{Conclusion}

\ifdefined\included
\else
\bibliographystyle{acm}
\bibliography{These}
\end{document}
\fi
